\section{Model and symmetries}
\label{sec:model}

We consider the dimerized XXZ chain with $L$ sites and open boundary
conditions, sites labeled $1,\dots,L$. The dimerized alternating-bond
structure goes back to Su, Schrieffer, and Heeger~\cite{su1979solitons}. With $X_m,Y_m,Z_m$ the Pauli
operators on site $m$, the Hamiltonian reads
\begin{equation}
H(s,\delta)=H_o+H_e,
\end{equation}
where $H_o$ ($H_e$) couples odd (even) bonds, $s$ is the dimerization
parameter, and $\delta$ controls the anisotropy:
\begin{equation}
H_o=(1-s)\sum_{j=1}^{L/2}h_{2j-1,2j},\qquad
H_e=s\sum_{j=1}^{L/2-1}h_{2j,2j+1},
\end{equation}
\begin{equation}
h_{kl}=e^{-\delta}(X_kX_l+Y_kY_l)+e^{+\delta}Z_kZ_l,
\end{equation}
where $h_{kl}$ is the two-site bond operator with
$J_{xx}=e^{-\delta}$ and $J_z=e^{+\delta}$; $\delta=0$ is the Heisenberg
isotropic point, while $\delta>0$ leans toward Ising.
Unless stated otherwise we take $L=8$.

$H$ commutes with the total magnetization
$Z_{\rm tot}=\sum_m Z_m$, the global spin flip $\prod_m X_m$, and the
full-chain reflection $R$ ($m\leftrightarrow L+1-m$), so evolution under
$H$ never leaves the initial symmetry sector. For $\delta\ne0$ the symmetry is $U(1)$ ($Z_{\rm tot}$
conserved); at $\delta=0$ it enlarges to $SU(2)$.

The topological phase here belongs to the one-dimensional
symmetry-protected family whose study began with Haldane's conjecture
of a finite gap for integer-spin antiferromagnetic
chains~\cite{haldane1983nonlinear,dennijs1989preroughening} and the
exact valence-bond-solid ground state found at a special point of the
spin-$1$ chain~\cite{affleck1987rigorous,dennijs1989preroughening}.
That phenomenology is now classified: gapped symmetric phases in one
dimension are labeled by the second cohomology group, i.e.\ projective
representations of the symmetry~\cite{chen2011classification}, equivalently
in matrix-product-state
form~\cite{schuch2011classifying}, with the degeneracy of the
entanglement spectrum as the operational
fingerprint~\cite{pollmann2010entanglement,li2008entanglement}.
Our dimerized XXZ chain realizes three competing orders---trivial,
topological, and antiferromagnetic---within a single tunable model,
which is why a marker that resolves all three at once is needed.

To distinguish the phases we use the normalized topological reflection
marker, following the partial-reflection many-body invariant of
Elben~\textit{et al.}~\cite{elben2020manybody}. Let $\rho_I$ be the reduced density matrix of the ground state
on the central block $I$ (for $L=8$, sites $3$--$6$) and
${\mathcal R}_I$ the mirror exchange about the central bond. With
$Z_{\mathcal R}={\rm Tr}(\rho_I{\mathcal R}_I)$ and the corresponding
quantities $\rho_{I_{1,2}}$ on the left/right halves of $I$,
\begin{equation}
\tilde{Z}_{\mathcal R}=
\frac{Z_{\mathcal R}}
{\sqrt{[{\rm Tr}(\rho_{I_1}^2)+{\rm Tr}(\rho_{I_2}^2)]/2}},
\end{equation}
where the denominator is the mean purity of the two halves.
In other words, $\tilde{Z}_{\mathcal R}$ is the reflection expectation
normalized by the local mixedness, so that it approaches $+1$ ($-1$) in
the trivial (topological) phase and $0$ in the magnetically ordered
phase, independent of the overall entanglement level.
The finite-size phase diagram mapped with this marker is reported in
Sec.~\ref{sec:phasediagram}.
